\section{布局与结构}

\begin{frame}{双栏布局}
  \begin{columns}[T]
    \begin{column}{0.48\textwidth}
      \begin{block}{左侧}
        适合放置方法概述、定义或图示说明。
      \end{block}
    \end{column}
    \begin{column}{0.48\textwidth}
      \begin{block}{右侧}
        适合放置实验现象、结论摘要或补充解释。
      \end{block}
    \end{column}
  \end{columns}
\end{frame}

\begin{frame}{三栏布局}
  \begin{columns}[T]
    \begin{column}{0.31\textwidth}
      \begin{exampleblock}{步骤一}
        输入与问题定义
      \end{exampleblock}
    \end{column}
    \begin{column}{0.31\textwidth}
      \begin{exampleblock}{步骤二}
        建模与求解过程
      \end{exampleblock}
    \end{column}
    \begin{column}{0.31\textwidth}
      \begin{exampleblock}{步骤三}
        输出与分析结论
      \end{exampleblock}
    \end{column}
  \end{columns}
\end{frame}

\begin{frame}{对比结构}
  \begin{columns}[T]
    \begin{column}{0.48\textwidth}
      \begin{alertblock}{传统写法}
        \begin{itemize}
          \item 内容和样式耦合
          \item 难以统一修改
          \item 长文档易失控
        \end{itemize}
      \end{alertblock}
    \end{column}
    \begin{column}{0.48\textwidth}
      \begin{exampleblock}{工程化写法}
        \begin{itemize}
          \item 主题统一控制样式
          \item 章节文件专注内容
          \item 更适合协作与复用
        \end{itemize}
      \end{exampleblock}
    \end{column}
  \end{columns}
\end{frame}
